\section{Estado, validacion y datos}

\subsection{Descripcion}

Este documento define como debe manejar el frontend el estado remoto, el estado local de interfaz, la validacion de entradas y el consumo de datos provenientes del backend. Su objetivo es evitar que el cliente mezcle preocupaciones, replique reglas oficiales de negocio o construya contratos ambiguos.

\subsection{Alcance}

Aplica al repositorio frontend en su integracion con backend. Describe principios y reglas de manejo de datos; no reemplaza la documentacion de contratos conceptuales del backend ni el detalle de implementacion de librerias.

\subsection{Definiciones}

\begin{itemize}
  \item \textbf{Estado remoto}: informacion obtenida del backend y considerada verdad vigente para el frontend.
  \item \textbf{Estado local}: estado efimero de interfaz usado para composicion visual, formularios y comportamiento interactivo no persistido.
  \item \textbf{Validacion de frontera}: validacion realizada en el punto de entrada o salida de datos entre frontend y backend.
  \item \textbf{Contrato}: estructura esperada de entrada o salida consumida por el frontend.
\end{itemize}

\subsection{Reglas}

\begin{itemize}
  \item El frontend debe tratar el estado remoto y el estado local como preocupaciones separadas.
  \item Los datos de dominio visibles en el cliente deben venir del backend y consumirse bajo contratos explicitamente validados.
  \item La validacion del frontend existe para mejorar experiencia y reducir errores de captura, no para sustituir enforcement del backend.
  \item El frontend no debe asumir como verdadero ningun lock, permiso, resultado o puntaje que no haya sido confirmado por el backend.
  \item Los errores funcionales y los errores tecnicos deben mostrarse de manera diferenciada en la experiencia de usuario.
\end{itemize}

\subsubsection{Manejo de estado remoto}

\begin{itemize}
  \item El estado remoto incluye, al menos, sesion, perfil, estado de activacion, calendario de partidos, picks propios, predicciones especiales, leaderboard y resultados visibles.
  \item El frontend debe considerar como fuente de verdad vigente el dato mas reciente validamente recibido del backend.
  \item Los datos remotos deben poder invalidarse o refrescarse cuando una mutacion cambie el estado funcional del sistema.
  \item Las vistas derivadas del frontend no deben persistirse como verdad paralela si pueden obtenerse de respuestas oficiales del backend.
\end{itemize}

\subsubsection{Manejo de estado local}

\begin{itemize}
  \item El estado local se limita a composicion visual, estado temporal de formularios, controles abiertos o cerrados, filtros no persistidos y feedback inmediato de interaccion.
  \item Un dato local no puede usarse como sustituto durable de una respuesta oficial del backend.
  \item Cuando una accion critica depende de informacion remota, el frontend debe reconciliar su estado local con la respuesta del servidor despues de la operacion.
\end{itemize}

\subsubsection{Reglas de validacion}

\begin{itemize}
  \item El frontend debe validar entradas visibles del usuario antes de enviar la solicitud cuando eso mejore claridad o reduzca errores evitables.
  \item La validacion de picks debe respetar que los goles permitidos son enteros entre \texttt{0} y \texttt{9}.
  \item La validacion de predicciones especiales debe respetar su cierre propio y el formato esperado de seleccion.
  \item La validacion de formularios no debe ocultar errores devueltos por backend; debe integrarlos y mostrarlos de forma comprensible.
  \item Si una operacion es rechazada por lock, permiso o estado funcional, el frontend debe privilegiar el mensaje funcional sobre el detalle tecnico.
\end{itemize}

\subsubsection{Contratos y acceso a datos}

\begin{itemize}
  \item El frontend consumira contratos conceptuales estables del backend y adaptara sus respuestas a modelos de presentacion cuando sea necesario.
  \item Las adaptaciones de contrato deben vivir cerca del acceso a datos o de la entidad correspondiente, no dispersas en componentes visuales.
  \item El frontend no debe construir contratos paralelos incompatibles con backend para simplificar localmente una pantalla.
  \item Toda mutacion relevante debe definir que consultas deben refrescarse o invalidarse despues de completarse.
\end{itemize}

\subsubsection{Datos minimos por feature}

\begin{longtable}{p{0.22\textwidth}p{0.32\textwidth}p{0.30\textwidth}}
\toprule
\textbf{Feature} & \textbf{Datos remotos principales} & \textbf{Estado local frecuente} \\
\midrule
\texttt{auth} & sesion y contexto basico de cuenta & estado de formulario y feedback de acceso \\
\texttt{activation} & estado de usuario, comprobante e instrucciones visibles & carga de archivo, confirmacion y mensajes de revision \\
\texttt{picks} & partidos, kickoff, picks propios y estado funcional del partido & valores editables, filtros y confirmaciones de guardado \\
\texttt{special-predictions} & predicciones especiales propias y su estado de cierre & seleccion temporal antes de guardar \\
\texttt{leaderboard} & ranking general, filtros y desglose disponible & filtros activos y preferencia de vista \\
\texttt{profile} & datos propios del participante & expansion de secciones y edicion local si aplica \\
\texttt{admin} & usuarios, comprobantes, partidos, resultados, recalculos y auditoria visible & filtros, confirmaciones y contexto de accion critica \\
\bottomrule
\end{longtable}

\subsubsection{Reglas de error y feedback}

\begin{itemize}
  \item Un error funcional debe explicar que regla impidio la accion, por ejemplo lock vencido, cuenta no activa o falta de permisos.
  \item Un error tecnico debe comunicar que la operacion no pudo completarse y ofrecer reintento cuando corresponda.
  \item El frontend no debe esconder el estado final de una mutacion critica; siempre debe dejar feedback de exito, rechazo o fallo tecnico.
  \item Cuando exista resultado corregido, el frontend debe poder mostrar un aviso explicito y breve sin reconstruir la auditoria interna.
\end{itemize}

\subsection{Edge Cases}

\begin{itemize}
  \item Si el frontend muestra un formulario editable y el backend rechaza por lock reciente, la vista debe reconciliarse con el nuevo estado funcional y pasar a modo lectura.
  \item Si un dato remoto cambia por recalculo o correccion, las vistas derivadas deben invalidarse para evitar mostrar ranking o puntajes obsoletos.
  \item Si una consulta administrativa retorna informacion parcial por permisos, el frontend debe degradar la vista de forma controlada y no asumir acceso total.
\end{itemize}

\subsection{Invariantes}

\begin{itemize}
  \item El backend sigue siendo la autoridad final para validar y persistir operaciones del dominio.
  \item Ningun componente del frontend debe depender de datos no garantizados por contrato.
  \item La logica de presentacion no debe contener reglas oficiales unicas de scoring, lock o permisos.
  \item Toda mutacion critica debe tener un camino claro de refresco o invalidacion del estado remoto afectado.
\end{itemize}

\subsection{Preguntas abiertas}

\begin{itemize}
  \item Confirmar si los esquemas Zod de frontend se mantendran totalmente independientes o si existira una estrategia posterior de comparticion controlada de contratos con backend.
  \item Confirmar cuantas superficies del frontend mostraran desglose de puntaje y con que precision visual se representaran los decimales.
\end{itemize}

\subsection{Decisiones pendientes}

\begin{itemize}
  \item \texttt{Decision pendiente} \texttt{No bloqueante}: definir la politica exacta de invalidacion y refresco para leaderboard y picks despues de una correccion de resultado.
  \item \texttt{Decision pendiente} \texttt{No bloqueante}: decidir si los mensajes funcionales del frontend tendran catalogo centralizado por feature o un inventario transversal de microcopy operativo.
\end{itemize}
